\paragraph{AQ-136-TENSOR: the flat tensor target already fails.}
The supplied SA08 intake already reports the corresponding
prime-power obstruction in its statement A-S09-26, at
\texttt{sections/A\_hadamard\_image.tex:25}. That prior full
proof is not present among the supplied sources. The proof here
rederives the tensor assertion completely and makes no novelty
claim; it does not certify the intake's additional assertion
about a full PD closure in a rational Hadamard algebra.

Let $k$ be a commutative ring, $S=\Gamma_k(T)$, and let $B$ be
the admissible divided-power simplex theory in AQ-136-DP-MODEL.
Form the actual simplicial PD differential $S$-algebra
\[
                         E=B\otimes_S B.
\]
Both inclusions $B\to E$ preserve the coefficient algebra.
Every simplex of $E$ is free over $S$ in each cochain degree,
and has vertical cohomology $S$ in degree zero and zero in
positive degrees, in every weight. Nevertheless, for every
rational prime $p$ and integer $r\geq1$,
\[
 H_1\bigl(N(E^{p^r,0}_*)\bigr)\simeq k/pk.                 \tag{136.T1}
\]
Consequently $E$ fails the horizontal admissibility condition
over every nonzero ring which is not a $\mathbb Q$-algebra.
In particular, this failure occurs even over the $p$-local
and positive-characteristic rings for which the original
common-source and terminal-groupoid questions are already
answered affirmatively above.

\begin{proof}
In simplex degree $n$, eliminate $x_0$ and $y_0$ using
$\sum x_j=\sum y_j=T$, and eliminate $e_0,f_0$ using
$\sum e_j=\sum f_j=0$. The ordinary tensor has the explicit
PD presentation
\[
 E_n=S\langle x_1,\ldots,x_n,y_1,\ldots,y_n\rangle
       \otimes_k\Lambda_k(e_1,\ldots,e_n,f_1,\ldots,f_n),
 \qquad dx_j=e_j,\quad dy_j=f_j.
\]
Here angle brackets mean free divided-power variables and
all the displayed variables have weight one. Their cochain
degrees are zero for $x_j,y_j$ and one for $e_j,f_j$.
This is a tensor of $2n$ divided-power de Rham disks over $S$.
Each such disk has augmentation to $S$, and its positive
part contracts by
\[
                    h(z^{[a]}\,dz)=z^{[a+1]}.
\]
This is an $S$-linear, weight-preserving contraction.
Tensoring the finite number of contractions proves the
asserted vertical cohomology and simplexwise freeness,
without a flatness hypothesis on $k$.

Put $P_*=B_*^{*,0}$, regarded now as a simplicial graded
$S$-module, and augment it to $k=S/(T^{[1]},T^{[2]},\ldots)$
by killing all positive weights. In simplex degree $n$,
\[
 P_n=\Gamma_k(x_0,\ldots,x_n)
       \simeq S\langle x_1,\ldots,x_n\rangle,
\]
so it is a free graded $S$-module. In weight $w>0$ its
underlying simplicial $k$-module is $\Gamma^w_k(L)$, with
$L_n=k^{n+1}$ and the summation simplicial maps already
specified in AQ-136-DP-MODEL. The explicit linear simplicial
contraction of $L$ induces a contraction of $\Gamma^w(L)$.
Weight zero is the constant $k$. Therefore $N(P_*)\to k$
is a resolution, and its terms are projective graded
$S$-modules: normalization is the image of a degreewise
splitting idempotent on $P_n$.

The zero-form row of $E$ is the diagonal tensor $P_*\otimes_S P_*$.
The Eilenberg--Zilber comparison for flat simplicial modules
identifies its normalized complex, up to quasi-isomorphism,
with
\[
                  N(P_*)\otimes_S N(P_*).
\]
An exact primary formulation of this input is Stacks,
Remark 21.39.11, Tag \texttt{08QD}, following
Lemma 21.39.10. The statement there is first made for the
associated unnormalized complexes; the split degenerate
summands give the normalized formulation used here.
The comparison respects the internal weight grading.
Both factors are nonnegative projective resolutions of
$k$, so their tensor computes $k\otimes_S^{\mathbf L}k$.
One may check the last assertion by the first-quadrant
double-complex filtration: tensoring one resolution with
each projective term of the other preserves its homology,
and each total degree has finitely many bidegrees. Thus
\[
 H_i\bigl(N(E^{*,0}_*)\bigr)
                         \simeq\operatorname{Tor}_i^S(k,k).
                                                               \tag{136.T2}
\]

Let $I=\bigoplus_{w>0}S^w$. Apply $-\otimes_S k$ to
$0\to I\to S\to k\to0$. Since $S$ is projective over
itself and the map $S\otimes_S k\to k\otimes_S k$ is
the identity of $k$, its exact Tor sequence gives
\[
                         \operatorname{Tor}_1^S(k,k)
                           \simeq I/I^2.
\]
For $m>1$, the weight-$m$ component of $I^2$ is generated
inside $kT^{[m]}$ by
\[
              \binom ma T^{[m]},\qquad 1\leq a<m.
\]
If $m=p^r$, every displayed binomial coefficient is divisible
by $p$, because $(1+X)^{p^r}=1+X^{p^r}$ in $\mathbb F_p[X]$.
Their integer greatest common divisor divides
$\binom{p^r}{1}=p^r$. Moreover
\[
 \binom{p^r}{p^{r-1}}
   =p\binom{p^r-1}{p^{r-1}-1}
\]
has $p$-adic valuation exactly one: modulo $p$, all the
coefficients of $(1+X)^{p^r-1}$ are alternately $1,-1$,
as follows by dividing $1+X^{p^r}$ by $1+X$.
The greatest common divisor is therefore exactly $p$.
Bezout's identity shows that these coefficients generate
the ideal $pk$ after any base change to $k$. Hence
$(I/I^2)^{p^r}=k/pk$, proving (136.T1).

If $k$ is a nonzero ring without a $\mathbb Q$-algebra
structure, some rational prime $p$ is not a unit in $k$;
otherwise the map from $\mathbb Z$ would extend to its
localization $\mathbb Q$. For that prime, $k/pk\ne0$.
Equation (136.T1) gives nonzero first simplicial homology
in arbitrarily high weights of the zero-form row. This
contradicts (136.4). \qedhere
\end{proof}

This is a failure of the displayed common-target construction,
including in a simplexwise free and vertically acyclic case.
Replacing a module tensor by its derived tensor does not
establish horizontal admissibility: (136.T2) already is the
derived-module calculation. No model structure on strict PD
cochain algebras, or claim that every possible PD derived
pushout presentation must have this zero-form row, is assumed.
In particular the result neither refutes the original
common-source conjecture nor the integral terminal-groupoid
assertion. Horizontal admissibility is additional information
about an actual simplicial presentation, not a consequence
of its simplexwise vertical derived isomorphism class.

